\documentclass[twoside,11pt]{article}

% Use JMLR style with preprint option
\usepackage[preprint]{jmlr2e}
\usepackage{amsmath}
\usepackage{bbm}
% Additional packages (jmlr2e already includes: epsfig, amssymb, natbib, graphicx)
\usepackage{algorithm}
\usepackage{algorithmic}
\usepackage{subcaption}
\usepackage{xcolor}
\usepackage{tikz}
\usetikzlibrary{calc, arrows.meta, decorations.markings, patterns, shapes.geometric, positioning, 3d, decorations.pathmorphing}
\usepackage{array}
\usepackage{tabularx}
\usepackage{listings}
\lstset{basicstyle=\ttfamily, breaklines=true}
\usepackage{booktabs}

% Custom commands (must match main paper)
\newcommand{\softmax}{\operatorname{softmax}}
\newcommand{\KL}{\operatorname{KL}}
\newcommand{\Tr}{\operatorname{Tr}}

\begin{document}

\jmlrheading{}{2026}{}{}{}{}{Robert C. Dennis}

\title{Attention as Gauge-Theoretic Variational Inference\\[6pt]
\large Supplementary Material}

\author{\name Robert C. Dennis \email cdenn016@gmail.com \\
       \addr Independent Researcher\\
       \addr Leander, Texas 78641, USA}

\editor{TBD}

\maketitle

\ShortHeadings{Gauge-Theoretic Attention: Supplementary Material}{Dennis}

This supplement collects derivations and implementation details that support the main paper but are not essential for following the core argument.  Appendix~A reviews standard differential geometry for readers from machine learning backgrounds. Appendix~B derives the covariance gradient and its equilibrium analysis.  Appendix~C derives the gauge frame gradients for both $\mathrm{SO}(N)$ and $\mathrm{GL}(K)$, including preconditioning strategies.  Appendix~D details the numerical methods used for variational gradient descent on the Gaussian and gauge-frame manifolds.

\appendix

\section{General Mathematical Framework}
\subsection{Principal Bundle and Associated Bundles}

The following geometric and probabilistic constructions comprise standard methods in differential geometry and gauge theory.  See references for details, proofs, and examples \citep{nakahara2003geometry,frankel2011geometry,blei2017variational,amari2016information}.

Let $\pi: \mathcal{N} \to \mathcal{C}$ be a smooth principal $G$-bundle where $\mathcal{C}$ is a smooth manifold (the base space) and $G$ is a Lie group (the structure group) acting freely and transitively on the right on $\mathcal{N}$. The projection satisfies $\pi(n \cdot g) = \pi(n)$ for all $g \in G$, $n \in \mathcal{N}$.

Let $\rho_q: G \to \text{Aut}(\mathcal{B}_q)$ and $\rho_p: G \to \text{Aut}(\mathcal{B}_p)$ be representations of $G$ on smooth statistical manifolds $\mathcal{B}_q$ (belief/recognition fiber) and $\mathcal{B}_p$ (model/prior fiber). These fibers are typically $K$-dimensional probability simplices $\Delta^K$ (for categorical distributions) or statistical manifolds with information-geometric structure (e.g., Gaussian manifolds, exponential families).

The associated bundles are:\\
\begin{align}
\mathcal{E}_q &:= \mathcal{N} \times_{\rho_q} \mathcal{B}_q = \left(\mathcal{N} \times \mathcal{B}_q\right)/\sim_q, \\
\mathcal{E}_p &:= \mathcal{N} \times_{\rho_p} \mathcal{B}_p = \left(\mathcal{N} \times \mathcal{B}_p\right)/\sim_p,
\end{align}\\
where $(n \cdot g, b) \sim_u (n, \rho_u(g)b)$ for $u \in \{q,p\}$.

These give fiber bundles $\pi_{\mathcal{E}_q}: \mathcal{E}_q \to \mathcal{C}$ and $\pi_{\mathcal{E}_p}: \mathcal{E}_p \to \mathcal{C}$ with fibers $\mathcal{B}_q(c) \cong \mathcal{B}_q$ and $\mathcal{B}_p(c) \cong \mathcal{B}_p$ at each $c \in \mathcal{C}$.

\subsection{Agents and Multi-Agent Systems}
Agent\\
An agent $\mathcal{A}^i$ is a pair of local sections over a domain $\mathcal{U}_i \subset \mathcal{C}$:

\begin{equation}
\mathcal{A}^i = \left(\sigma^i_q, \sigma^i_p\right),
\end{equation}

where $\sigma^i_q : \mathcal{U}_i \to \mathcal{E}_q$ and $\sigma^i_p : \mathcal{U}_i \to \mathcal{E}_p$.

We write $q_i(c) := \sigma^i_q(c) \in \mathcal{B}_q(c)$ and $p_i(c) := \sigma^i_p(c) \in \mathcal{B}_p(c)$ for the belief and model at base point $c$.

Multi-Agent System\\
A multi-agent system $\mathcal{M}$ over $\mathcal{C}$ is a collection of agents indexed by $\mathcal{I}$:

\begin{equation}
\mathcal{M} = \left\{\mathcal{A}^i = \left(\sigma^i_q, \sigma^i_p\right)\right\}_{i \in \mathcal{I}}.
\end{equation}

Agents generally overlap on intersections $\mathcal{U}_i \cap \mathcal{U}_j$.

Meta-Agent and Epistemic Death\\
A meta-agent is a multi-agent system whose component agents share identical section values on their overlap:

\begin{equation}
q_i(c) = q_j(c), \quad p_i(c) = p_j(c) \quad \text{for } c \in \mathcal{U}_i \cap \mathcal{U}_j.
\end{equation}

A set of agents is epistemically dead if they identically share both beliefs and models. While constituent agents of a meta-agent may be epistemically dead, the meta-agent itself need not be; such agents can be integrated out, yielding coarse-grained higher-order entities and a route toward emergence.

\subsubsection{Hierarchical Meta-Agent Emergence and Cross-Scale Coupling}
\label{sec:appendix_meta_agent}

The gauge-theoretic framework naturally supports the emergence of meta-agents through coarse-graining \citep{shen2008coarse}. A meta-agent $\mathcal{M}^{(1)}$ is a set of agents $\{i \in I_{\mathcal{M}}\}$ satisfying:

\begin{align}
q_i &= \Omega_{ij} q_j \quad \forall i,j \in I_{\mathcal{M}} \quad \text{(belief consensus)}, \label{eq:belief_consensus} \\
s_i &= \tilde{\Omega}_{ij} s_j \quad \forall i,j \in I_{\mathcal{M}} \quad \text{(model consensus)}. \label{eq:model_consensus}
\end{align}

When these conditions hold, the agents are epistemically dead in that they share identical beliefs and models after accounting for gauge transformations and no longer evolve without continual observations. The meta-agent can be described by a single representative renormalized state $(q_{\mathcal{M}}, s_{\mathcal{M}})$.

Meta-agents at scale $\zeta=0$ (individual agents) can form meta-agents at scale $\zeta=1$ (groups), which can further coalesce into meta-agents at scale $\zeta=2$ (communities), and so on, creating a hierarchical scale structure:

\begin{figure}[h]
\centering
\small
\[
\text{Agents}^{(0)} \xrightarrow{\text{consensus}}
\text{Meta-agents}^{(1)} \xrightarrow{\text{consensus}}
\text{Meta}^{(2)} \xrightarrow{\text{consensus}} \cdots
\]
\caption{Hierarchical emergence through consensus at successive scales $\zeta$.}
\label{fig:hierarchy}
\end{figure}

At each scale, the effective free energy takes the same form as the full variational free energy (Eq.~7 of the main paper), with agents at scale $\zeta$ replaced by meta-agents at scale $\zeta+1$, coupling constants renormalized according to
\begin{equation*}
  \beta_{ij}^{(\zeta+1)} = f_\beta(\{\beta_{kl}^{(\zeta)}\}), \quad
  \gamma_{ij}^{(\zeta+1)} = f_\gamma(\{\gamma_{kl}^{(\zeta)}\}),
\end{equation*}
and effective gauge frames given by $\phi_{\mathcal{M}}^{(\zeta+1)} = \text{average}(\{\phi_i^{(\zeta)}\})$.
This is the gauge-theoretic analogue of renormalization group flow in statistical field theory \citep{anderson1984basic,wilson1974renormalization,garciaMillan2024network}. Cross-scale couplings $\Lambda^{\zeta}_{\zeta'}$ mediate interactions between agents at different hierarchical levels via the cross-scale transport operators defined below.

Whether this hierarchical structure maps onto multi-layer transformer architectures (where each layer constitutes a ``scale''), or onto distinct organizational levels within a single attention layer, remains an open question for future investigation.

\subsection{Bundle Morphisms and Transport Operators}
Via standard horizontal lifting from the principal bundle $\mathcal{N}$ to the associated bundles, we obtain a hierarchy of morphisms and induced connections:

Intra-bundle transport operators $\Omega^{(q)}_{ij}: \Gamma(\mathcal{B}_q) \to \Gamma(\mathcal{B}_q)$ (belief-to-belief) and $\Omega^{(p)}_{ij}: \Gamma(\mathcal{B}_p) \to \Gamma(\mathcal{B}_p)$ (model-to-model) act within a single bundle. Cross-scale transport operators $\Lambda^{s}_{s'}: \Gamma^{s}(\mathcal{B}_q) \to \Gamma^{s'}(\mathcal{B}_q)$ (beliefs across scales) and $\tilde{\Lambda}^{s}_{s'}: \Gamma^{s}(\mathcal{B}_p) \to \Gamma^{s'}(\mathcal{B}_p)$ (models across scales) connect different hierarchical levels. Inter-bundle morphisms $\Theta^i_j: \Gamma(\mathcal{B}_q) \to \Gamma(\mathcal{B}_p)$ (belief to model) and $\tilde{\Theta}^i_j: \Gamma(\mathcal{B}_p) \to \Gamma(\mathcal{B}_q)$ (model to belief) relate the two fiber types. Finally, global bundle morphisms $\Phi: \mathcal{E}_p \to \mathcal{E}_q$ (model bundle to belief bundle) and $\tilde{\Phi}: \mathcal{E}_q \to \mathcal{E}_p$ (belief bundle to model bundle) act on the total associated bundles.

Here $\Gamma(\mathcal{B}_u)$ denotes the space of smooth sections over the base $\mathcal{C}$.

\subsection{Gauge Frames and Connections}
Each agent $i$ possesses a local gauge frame field:

\begin{equation}
\phi_i: \mathcal{U}_i \to \mathfrak{g} = \text{Lie}(G),
\end{equation}

which induces a local connection one-form:

\begin{equation}
A^{(i)}_\mu(c) = U_i^{-1}(c) \partial_\mu U_i(c),
\end{equation}

where $U_i(c) = \exp[\phi_i(c)] \in G$. The associated field strength (gauge curvature) is:

\begin{equation}
F^{(i)}_{\mu\nu}(c) = \partial_\mu A^{(i)}_\nu - \partial_\nu A^{(i)}_\mu + [A^{(i)}_\mu, A^{(i)}_\nu] \in \mathfrak{g}.
\end{equation}

When two agents overlap at $c \in \mathcal{U}_i \cap \mathcal{U}_j$, the inter-agent gauge transformation:

\begin{equation}
\Omega_{ij}(c) = \exp[\phi_i(c)] \exp[-\phi_j(c)] \in G
\end{equation}

transports agent $j$'s representations into agent $i$'s frame:

\begin{equation}
q_j(c) \mapsto \Omega_{ij}(c) \cdot q_j(c) := \rho(\Omega_{ij}(c)) q_j(c).
\end{equation}

On overlaps, local connections are related by:

\begin{equation}
A^{(i)}_\mu = \Omega_{ij} A^{(j)}_\mu \Omega_{ij}^{-1} + \Omega_{ij} \partial_\mu \Omega_{ij}^{-1}.
\end{equation}



\section{Covariance Dynamics and Equilibrium Analysis}
\label{app:covariance_dynamics}

\subsection{Covariance Gradient of the Generalized Free Energy}

Here we derive the gradient of the single-agent free energy $\mathcal{F}_i$ with respect to the covariance $\Sigma_i$ for agent $i$; a well known result in information geometry.

The free energy decomposes as

\begin{equation}
\mathcal{F}_i
=D_{\mathrm{KL}}(q_i \,\|\, p_i)
+ \sum_{j \neq i} \beta_{ij}\,
D_{\mathrm{KL}}(q_i \,\|\, \Omega_{ij} q_j)
- \mathbb{E}_{q_i}[\log p(o_i \mid k_i)],
\end{equation}

where $q_i = \mathcal{N}(\mu_i, \Sigma_i)$, $p_i = \mathcal{N}(\mu_{p,i}, \Sigma_{p,i})$, and $\sum_j \beta_{ij} = 1$ by construction.

\subsubsection{Gaussian KL Divergence under $\mathrm{GL}(K)$ Gauge Transport}

The alignment term $D_{\mathrm{KL}}(q_i \,\|\, \Omega_{ij} q_j)$ involves KL divergences between Gaussians related by gauge transport $\Omega_{ij} = e^{\phi_i}e^{-\phi_j} \in \mathrm{GL}^+(K)$. The standard Gaussian KL divergence is

\begin{align}
D_{\mathrm{KL}}\!\left(
\mathcal{N}(\mu_1, \Sigma_1)
\,\|\,
\mathcal{N}(\mu_2, \Sigma_2)
\right)
&=
\frac{1}{2}
\Big[
\log \frac{|\Sigma_2|}{|\Sigma_1|}
+ \mathrm{tr}\!\big(\Sigma_2^{-1}\Sigma_1\big)
\\ &\quad
+ (\mu_2 - \mu_1)^\top
  \Sigma_2^{-1}
  (\mu_2 - \mu_1)
- d
\Big].
\label{eq:gaussian_KL}
\end{align}


and differentiating w.r.t.\ $\Sigma_1$ (holding $\mu_1,\mu_2,\Sigma_2$ fixed) gives

\begin{equation}
\frac{\partial D_{\mathrm{KL}}}{\partial \Sigma_1}
=\frac{1}{2}\left[
-\Sigma_1^{-1}
+ \Sigma_2^{-1}
\right].
\end{equation}

Applying this to each term in $\mathcal{F}_i$, where the alignment KL involves the $\mathrm{GL}(K)$-transported belief $\Omega_{ij} q_j = \mathcal{N}(\Omega_{ij}\mu_j, \Omega_{ij}\Sigma_j\Omega_{ij}^\top)$:

\begin{align}
\frac{\partial}{\partial \Sigma_i}
D_{\mathrm{KL}}(q_i \,\|\, p_i)
&=\frac{1}{2}\left[
-\Sigma_i^{-1}
+ \Sigma_{p,i}^{-1}
\right],
\\[6pt]
\frac{\partial}{\partial \Sigma_i}
D_{\mathrm{KL}}(q_i \,\|\, \Omega_{ij} q_j)
&=
\frac{1}{2}\left[
-\Sigma_i^{-1}
+ (\Omega_{ij}\Sigma_j\Omega_{ij}^\top)^{-1}
\right].
\end{align}

The observation term contributes an $O(\Sigma_i^{-1})$ correction that we neglect in the high-precision / strong-alignment regime. The alignment term $\sum_j \beta_{ij} D_{\mathrm{KL}}(q_i \,\|\, \Omega_{ij} q_j)$ requires the product rule, since $\beta_{ij} = \mathrm{softmax}_j(-D_{\mathrm{KL}}/\tau)$ depends on $\Sigma_i$ through the KL divergence. Thus

\begin{equation}
\boxed{
\begin{aligned}
\frac{\partial \mathcal{F}_i}{\partial \Sigma_i}
&=\frac{1}{2}
\left[
-2 \Sigma_i^{-1}
+ \Sigma_{p,i}^{-1}
+ \sum_j \beta_{ij}
(\Omega_{ij}\Sigma_j\Omega_{ij}^\top)^{-1}
\right] \\
&\quad + \sum_j \frac{\partial \beta_{ij}}{\partial \Sigma_i}
D_{\mathrm{KL}}(q_i \,\|\, \Omega_{ij} q_j).
\end{aligned}
}
\label{eq:Sigma_gradient_final}
\end{equation}

The first line contains the ``direct'' gradient terms (differentiating each $K_{ij}$ at fixed $\beta_{ij}$). The second line is the attention weight gradient, analogous to the $\partial \beta_{ij}/\partial \mu_i$ term in the mean gradient. From the softmax structure, $\sum_j \partial \beta_{ij}/\partial \Sigma_i = 0$, so this correction vanishes when the KL divergences are uniform across neighbors, and also in both the $\tau \to 0$ and $\tau \to \infty$ limits.

Because $\sum_j \beta_{ij} = 1$, there is no $(1 + \sum_j \beta_{ij})$ prefactor in front of $\Sigma_i^{-1}$.

\subsection{Fixed-Point Equation and Symmetric Solution}

At equilibrium, we set $\partial \mathcal{F}_i / \partial \Sigma_i = 0$. Since $\sum_j \partial \beta_{ij}/\partial \Sigma_i = 0$, the attention weight correction term vanishes whenever the KL divergences $D_{\mathrm{KL}}(q_i \,\|\, \Omega_{ij}q_j)$ are approximately uniform across neighbors, and also in both the hard-attention ($\tau \to 0$) and uniform-attention ($\tau \to \infty$) limits. In these regimes, the equilibrium condition reduces to

\begin{equation}
\Sigma_i^{-1}
=\frac{1}{2}
\left[
\Sigma_{p,i}^{-1}
+ \sum_j \beta_{ij}
(\Omega_{ij}\Sigma_j\Omega_{ij}^\top)^{-1}
\right].
\label{eq:sigma_fixed_point_beta}
\end{equation}

This is a matrix-valued fixed-point equation coupling all agents: each agent's precision is the $\beta$-weighted combination of its own prior precision and the transported neighbor precisions. At intermediate temperatures, the attention weight correction provides an additional contribution that couples precision dynamics to the variance of KL divergences across neighbors.

\subsubsection{Homogeneous limit}

In the homogenous limit we assume

(i) all agents are identical, so $\Sigma_i = \Sigma_\infty$ for all $i$,

(ii) $\Omega_{ij} \approx I$ (weak misalignment),

(iii) $\Sigma_{p,i} = \Sigma_0$ (shared prior).

Then \eqref{eq:sigma_fixed_point_beta} becomes

\begin{equation}
\Sigma_\infty^{-1}
=\frac{1}{2}
\left[
\Sigma_0^{-1}
+ \Sigma_\infty^{-1}
\right]
\quad\Longrightarrow\quad
\Sigma_\infty = \Sigma_0.
\end{equation}

Hence, in a perfectly symmetric population, the equilibrium covariance reproduces the shared prior.

\subsubsection{Alignment-dominated regime}

Although $\sum_j \beta_{ij} = 1$, the effective strength of alignment is controlled by the parameter $\tau$ in

\begin{equation}
\beta_{ij}
=\frac{
\exp\!\left[
-\frac{1}{\tau}
D_{\mathrm{KL}}\!\big(q_i \,\|\, \Omega_{ij} q_j\big)
\right]
}{\sum_k
\exp\!\left[
-\frac{1}{\tau}
D_{\mathrm{KL}}\!\big(q_i \,\|\, \Omega_{ik} q_k\big)
\right]}.
\end{equation}

As $\tau \to 0$, $\beta_{ij}$ becomes sharply peaked on whichever neighbor $j$ best agrees (after transport). In that low-$\tau$ limit, the prior precision $\Sigma_{p,i}^{-1}$ becomes negligible relative to the socially enforced alignment term, and

\begin{equation}
\boxed{
\Sigma_i^{-1}
\;\approx\;
\sum_j \beta_{ij}
(\Omega_{ij}\Sigma_j\Omega_{ij}^\top)^{-1}
\;=\;\big\langle(\Omega_{ij}\Sigma_j\Omega_{ij}^\top)^{-1}
\big\rangle_{\beta}.}
\label{eq:beta_weighted_precision}
\end{equation}

Here $\langle \cdot \rangle_{\beta}$ denotes a $\beta$-weighted expectation over neighbors.

Therefore, in the strong-alignment (i.e. small $\tau$) regime, agent $i$'s precision matrix becomes the $\beta$-weighted average of its neighbors' transported precisions.

When, additionally, all agents already have approximately equal covariances and the transports are near-identity, i.e. $\Sigma_i \approx \Sigma_j$ and $\Omega_{ij} \approx I$, then \eqref{eq:beta_weighted_precision} implies

\begin{equation}
\Sigma_i^{-1} \approx \Sigma_j^{-1}
\quad\Longrightarrow\quad
\Sigma_i \approx \Omega_{ij}\Sigma_j\Omega_{ij}^\top,
\end{equation}

justifying the alignment assumption used in the main text.

\subsubsection{Gradient flow dynamics}

Finally, consider the gradient flow

\begin{equation}
\frac{d\Sigma_i}{dt}
=-\eta_\Sigma
\frac{\partial \mathcal{F}_i}{\partial \Sigma_i},
\qquad\eta_\Sigma > 0.
\end{equation}

Local stability of the equilibrium \eqref{eq:sigma_fixed_point_beta} follows from the positive-definiteness of the Hessian. For the Gaussian KL terms,

\begin{equation}
\frac{\partial^2 D_{\mathrm{KL}}}{\partial \Sigma_1 \, \partial \Sigma_1}
\;\sim\;
\Sigma_1^{-1} \otimes \Sigma_1^{-1}
+ \Sigma_2^{-1} \otimes \Sigma_2^{-1},
\end{equation}

which is manifestly positive definite for $\Sigma_1,\Sigma_2 \succ 0$.

Hence the covariance alignment fixed-point is an attractor of the variational dynamics.

Therefore, we find that $\Sigma_i \approx \Omega_{ij}\Sigma_j\Omega_{ij}^\top$ emerges from the dynamics itself rather than as an imposed constraint.


\section{Gauge Frame Gradients}
\label{app:phi_gradients}

We derive the gradient of the variational free energy $\mathcal{F}$ with respect to the gauge frame parameters $\phi_i \in \mathfrak{g}$. Unlike the covariance gradients (Appendix~B), gauge frame gradients require the differential of the matrix exponential map, and the resulting expressions differ for compact (e.g., $\mathrm{SO}(N)$) and general ($\mathrm{GL}(K)$) gauge groups.

\subsection{Structure of the Gauge Frame Gradient}

The gauge frames enter the free energy exclusively through the parallel transport operators $\Omega_{ij} = \exp(\phi_i)\exp(-\phi_j)$. Since neither the belief-prior term $D_{\mathrm{KL}}(q_i \| p_i)$ nor the observation likelihood $-\mathbb{E}_{q_i}[\log p(o_i | k_i)]$ depends on $\phi_i$, the complete gradient is:

\begin{equation}
\frac{\partial \mathcal{F}}{\partial \phi_i}
= \sum_j \left[
\frac{\partial \beta_{ij}}{\partial \phi_i} K_{ij}
+ \beta_{ij} \frac{\partial K_{ij}}{\partial \phi_i}
\right]
+ \sum_k \left[
\beta_{ki} \frac{\partial K_{ki}}{\partial \phi_i}
+ \sum_l \frac{\partial \beta_{kl}}{\partial \phi_i} K_{kl}
\right],
\label{eq:phi_grad_complete}
\end{equation}

where $K_{ij} \equiv D_{\mathrm{KL}}(q_i \| \Omega_{ij} q_j)$. The first sum captures agent $i$ aligning to others, with the product rule applied to $\beta_{ij} K_{ij}$. The second sum captures others aligning to $i$: changing $\phi_i$ shifts $K_{ki}$ through $\Omega_{ki} = \exp(\phi_k)\exp(-\phi_i)$, which both directly adjusts the alignment cost (the $\beta_{ki}\,\partial K_{ki}/\partial \phi_i$ term) and redistributes all of agent $k$'s attention weights (the $\sum_l (\partial \beta_{kl}/\partial \phi_i)\, K_{kl}$ term). Since only $K_{ki}$ depends on $\phi_i$ in agent $k$'s softmax, the latter simplifies to:

\begin{equation}
\sum_l \frac{\partial \beta_{kl}}{\partial \phi_i} K_{kl}
= -\frac{\beta_{ki}}{\tau}\frac{\partial K_{ki}}{\partial \phi_i}
\left(K_{ki} - \sum_l \beta_{kl} K_{kl}\right).
\label{eq:reverse_beta_grad_phi}
\end{equation}

This correction vanishes when $K_{ki}$ equals agent $k$'s average divergence $\bar{K}_k = \sum_l \beta_{kl} K_{kl}$, or when $\beta_{ki} \to 0$ (agent $k$ does not attend to $i$).

The coupling weight gradient in the first sum follows from the softmax structure:

\begin{equation}
\frac{\partial \beta_{ij}}{\partial \phi_i} = -\frac{\beta_{ij}}{\tau}\left[
\frac{\partial K_{ij}}{\partial \phi_i} - \sum_k \beta_{ik} \frac{\partial K_{ik}}{\partial \phi_i}
\right].
\label{eq:beta_grad_phi}
\end{equation}

\subsection{Differential of the Matrix Exponential}

All KL gradient terms require the derivative of the transport operator with respect to the gauge frame. For $\Omega_{ij} = \exp(\phi_i)\exp(-\phi_j)$:

\begin{align}
\frac{\partial \Omega_{ij}}{\partial \phi_i} &= \mathrm{dexp}_{\phi_i}(\cdot) \cdot \exp(-\phi_j), \label{eq:dOmega_dphi_i} \\
\frac{\partial \Omega_{ij}}{\partial \phi_j} &= -\exp(\phi_i) \cdot \mathrm{dexp}_{\phi_j}(\cdot) \cdot \exp(-\phi_j), \label{eq:dOmega_dphi_j}
\end{align}

where $\mathrm{dexp}_\phi: \mathfrak{g} \to \mathfrak{g}$ denotes the differential of the exponential map at $\phi$. For a direction $\xi \in \mathfrak{g}$, the differential is given by \citep{gallier2020differential,rossmann2002lie}:

\begin{equation}
\mathrm{dexp}_\phi(\xi) = \sum_{n=0}^{\infty} \frac{1}{(n+1)!} \mathrm{ad}_\phi^n(\xi) = \frac{e^{\mathrm{ad}_\phi} - I}{\mathrm{ad}_\phi}(\xi),
\label{eq:dexp_series}
\end{equation}

where $\mathrm{ad}_\phi(\xi) = [\phi, \xi] = \phi\xi - \xi\phi$ is the adjoint representation.

\subsubsection{$\mathrm{SO}(N)$ specialization.}

For $\mathrm{SO}(N)$, the gauge frame is expanded in skew-symmetric generators: $\phi = \sum_a \phi^a T_a$ where $T_a \in \mathfrak{so}(N)$, and the gradient is computed component-wise as $\partial \mathcal{F}/\partial \phi^a$. For $\mathrm{SO}(3)$ with $\|\phi\| = \theta$, the differential admits a closed form via the Rodrigues formula:

\begin{equation}
\mathrm{dexp}_\phi(T_a) = T_a + c_1(\theta)[\phi, T_a] + c_2(\theta)[\phi, [\phi, T_a]],
\label{eq:dexp_so3}
\end{equation}

with coefficients:

\begin{equation}
c_1(\theta) = \frac{1 - \cos\theta}{\theta^2}, \qquad c_2(\theta) = \frac{\theta - \sin\theta}{\theta^3}.
\end{equation}

For $\theta < \epsilon$ (numerically $\epsilon \sim 10^{-4}$), Taylor expansions $c_1 \approx \frac{1}{2} - \frac{\theta^2}{24}$ and $c_2 \approx \frac{1}{6} - \frac{\theta^2}{120}$ provide stable evaluation.

\subsubsection{$\mathrm{GL}(K)$ specialization.}

For $\mathrm{GL}(K)$, the gauge frame $\phi \in \mathfrak{gl}(K) \cong \mathbb{R}^{K \times K}$ is a general $K \times K$ matrix. The differential of the matrix exponential does not simplify to a Rodrigues-like formula, and is instead computed via the integral representation:

\begin{equation}
\mathrm{dexp}_\phi(\xi) = \int_0^1 e^{t\phi} \, \xi \, e^{(1-t)\phi} \, dt,
\label{eq:dexp_integral}
\end{equation}

or equivalently through the $(2K \times 2K)$ block matrix identity \citep{higham2008functions}:

\begin{equation}
\exp\!\begin{pmatrix} \phi & \xi \\ 0 & \phi \end{pmatrix}
= \begin{pmatrix} e^\phi & \mathrm{dexp}_\phi(\xi) \\ 0 & e^\phi \end{pmatrix}.
\label{eq:dexp_block}
\end{equation}

In practice, automatic differentiation through \texttt{torch.matrix\_exp} provides the $\mathrm{GL}(K)$ differential without explicit computation of \eqref{eq:dexp_integral} or \eqref{eq:dexp_block}.

\subsection{KL Gradient Through Transport}

The KL divergence $K_{ij} = D_{\mathrm{KL}}(q_i \| \Omega_{ij} q_j)$ between $q_i = \mathcal{N}(\mu_i, \Sigma_i)$ and the transported belief $\Omega_{ij} q_j = \mathcal{N}(\Omega_{ij}\mu_j, \Omega_{ij}\Sigma_j\Omega_{ij}^\top)$ depends on $\phi_i$ through $\Omega_{ij}$. Using the Gaussian KL \eqref{eq:gaussian_KL} and the chain rule:

\begin{align}
\frac{\partial K_{ij}}{\partial \phi_i^a}
&= \frac{\partial}{\partial \phi_i^a}\left\{
\frac{1}{2}\left[
\log\frac{|\Omega_{ij}\Sigma_j\Omega_{ij}^\top|}{|\Sigma_i|}
+ \mathrm{tr}\!\left((\Omega_{ij}\Sigma_j\Omega_{ij}^\top)^{-1}\Sigma_i\right)
+ \Delta\mu_{ij}^\top (\Omega_{ij}\Sigma_j\Omega_{ij}^\top)^{-1} \Delta\mu_{ij}
- K
\right]\right\},
\label{eq:dKL_dphi_expanded}
\end{align}

where $\Delta\mu_{ij} = \mu_i - \Omega_{ij}\mu_j$. Let $\tilde{\Sigma}_{ij} \equiv \Omega_{ij}\Sigma_j\Omega_{ij}^\top$ and $\tilde{\mu}_{ij} \equiv \Omega_{ij}\mu_j$ denote the transported statistics. The derivative has three contributions:

\paragraph{Mean term.} From the Mahalanobis distance:

\begin{equation}
\frac{\partial}{\partial \phi_i^a}\left[(\mu_i - \tilde{\mu}_{ij})^\top \tilde{\Sigma}_{ij}^{-1}(\mu_i - \tilde{\mu}_{ij})\right]
= -2(\mu_i - \tilde{\mu}_{ij})^\top \tilde{\Sigma}_{ij}^{-1} \frac{\partial \tilde{\mu}_{ij}}{\partial \phi_i^a}
+ (\text{terms from } \partial\tilde{\Sigma}_{ij}^{-1}/\partial\phi_i^a).
\label{eq:mean_term_phi}
\end{equation}

\paragraph{Trace term.} From the precision-covariance product:

\begin{equation}
\frac{\partial}{\partial \phi_i^a}\mathrm{tr}(\tilde{\Sigma}_{ij}^{-1}\Sigma_i)
= -\mathrm{tr}\!\left(\tilde{\Sigma}_{ij}^{-1} \frac{\partial \tilde{\Sigma}_{ij}}{\partial \phi_i^a} \tilde{\Sigma}_{ij}^{-1} \Sigma_i\right).
\label{eq:trace_term_phi}
\end{equation}

\paragraph{Log-determinant term.} From the log-determinant ratio:

\begin{equation}
\frac{\partial}{\partial \phi_i^a}\log|\tilde{\Sigma}_{ij}|
= \mathrm{tr}\!\left(\tilde{\Sigma}_{ij}^{-1} \frac{\partial \tilde{\Sigma}_{ij}}{\partial \phi_i^a}\right).
\label{eq:logdet_term_phi}
\end{equation}

\paragraph{Transport derivatives.} The derivatives of the transported statistics with respect to $\phi_i^a$ are:

\begin{align}
\frac{\partial \tilde{\mu}_{ij}}{\partial \phi_i^a} &= \frac{\partial (\Omega_{ij}\mu_j)}{\partial \phi_i^a}
= \mathrm{dexp}_{\phi_i}(T_a) \cdot \exp(-\phi_j) \, \mu_j
\equiv Q_a^{(i)} \, \exp(-\phi_j) \, \mu_j, \label{eq:dtilde_mu} \\
\frac{\partial \tilde{\Sigma}_{ij}}{\partial \phi_i^a} &=
Q_a^{(i)} \exp(-\phi_j) \Sigma_j \exp(-\phi_j)^\top \Omega_{ij}^\top
+ \Omega_{ij} \exp(-\phi_j) \Sigma_j \exp(-\phi_j)^\top (Q_a^{(i)})^\top, \label{eq:dtilde_Sigma}
\end{align}

where $Q_a^{(i)} \equiv \mathrm{dexp}_{\phi_i}(T_a)$ is the differential of the exponential map evaluated at $\phi_i$ in the direction of generator $T_a$.

\subsection{Retraction and Numerical Considerations}

\paragraph{Lie algebra updates.} Since $\phi_i \in \mathfrak{g}$ (a vector space), gradient descent proceeds in the Lie algebra directly:

\begin{equation}
\phi_i \leftarrow \phi_i - \eta_\phi \frac{\partial \mathcal{F}}{\partial \phi_i},
\label{eq:phi_update}
\end{equation}

where $\eta_\phi$ is the gauge frame learning rate. No metric correction (Fisher-Rao or otherwise) is applied: the exponential map $\exp: \mathfrak{g} \to G$ provides natural coordinates for the Lie group, and the gradient naturally lives in $\mathfrak{g}$.

\paragraph{$\mathrm{SO}(N)$ retraction.} For compact groups such as $\mathrm{SO}(N)$, the exponential map is periodic. To ensure unique parameterization and avoid branch cuts, we retract to the principal ball $\|\phi\| < \pi - \epsilon$ (with $\epsilon \sim 10^{-2}$) after each update. If $\|\phi\| \geq \pi$, the angle is wrapped modulo $2\pi$ and, if necessary, the axis is flipped via antipodal identification ($\phi \to -\phi$ when $\theta > \pi$), exploiting the fact that $\exp(\phi) = \exp(\phi - 2\pi \hat{n})$ for $\mathrm{SO}(3)$.

\paragraph{$\mathrm{GL}(K)$ updates.} For $\mathrm{GL}(K)$, the matrix exponential is surjective onto $\mathrm{GL}^+(K)$ (the identity component), and no retraction is required beyond standard gradient clipping for numerical stability. In our implementation, automatic differentiation through PyTorch's \texttt{matrix\_exp} handles the $\mathrm{GL}(K)$ differential implicitly.

\subsection{Gauge Frame Preconditioning for $\mathrm{GL}(K)$}
\label{sec:glk_preconditioning}

While the basic update~\eqref{eq:phi_update} uses the Euclidean gradient in the Lie algebra, the non-compact structure of $\mathrm{GL}(K)$ introduces a numerical asymmetry that motivates gauge-aware preconditioning. The Lie algebra decomposes via the Cartan involution $\theta(X) = -X^\top$ as:

\begin{equation}
\mathfrak{gl}(K) = \underbrace{\mathfrak{so}(K)}_{\text{compact (skew-symmetric)}} \;\oplus\; \underbrace{\mathrm{Sym}(K)}_{\text{non-compact (symmetric)}} \;\oplus\; \underbrace{\mathbb{R}}_{\text{center (trace)}},
\label{eq:cartan_decomposition}
\end{equation}

where $\mathfrak{so}(K) = \{X : X + X^\top = 0\}$ generates rotations and $\mathrm{Sym}(K) = \{X : X = X^\top, \operatorname{tr} X = 0\}$ generates volume-preserving shears. These two subspaces have qualitatively different gradient geometry: for the compact directions, $\|d\exp_\phi(T_a)\| = O(1)$ uniformly in $\|\phi\|$, while for the symmetric directions, $\|d\exp_\phi(T_a)\| \sim \exp(\|\phi\|)$. Euclidean gradient descent therefore systematically over-steps in the non-compact directions relative to the compact ones, leading to numerical instability for $\mathrm{GL}(K)$ gauge frames.

We implement four preconditioning strategies of increasing geometric fidelity:

\paragraph{1. Norm clipping (baseline).}
Simple gradient norm clipping $\tilde{g} = g \cdot \min(1, c/\|g\|)$ with threshold $c$. No geometric awareness; treats all Lie algebra directions uniformly.

\paragraph{2. Cartan preconditioning.}
Decompose the gradient into compact and non-compact components and selectively dampen the latter. Given Lie algebra generators $\{T_a\}_{a=1}^{n_{\mathrm{gen}}}$ with Gram matrix $G_{ab} = \operatorname{tr}(T_a^\top T_b)$ and symmetric inner product $S_{ab} = \operatorname{tr}(T_a T_b)$, the projector onto the symmetric subspace is:

\begin{equation}
P_{\mathrm{sym}} = \tfrac{1}{2}G^{-1}(G + S),
\end{equation}

and the preconditioner is $M = I - (1 - \lambda_{\mathrm{sym}}) P_{\mathrm{sym}}$ with dampening factor $\lambda_{\mathrm{sym}} \in (0, 1)$ (typically $0.1$, corresponding to $10\times$ dampening of symmetric directions). The preconditioned gradient is $\tilde{\nabla}_\phi \mathcal{F} = M \cdot \nabla_\phi \mathcal{F}$.

\paragraph{3. Killing form natural gradient.}
The Killing form of $\mathfrak{gl}(K)$ is $\kappa(X, Y) = 2K \operatorname{tr}(XY) - 2\operatorname{tr}(X)\operatorname{tr}(Y)$. In the generator basis, this yields the metric:

\begin{equation}
\tilde{g}_{ab} = 2K \operatorname{tr}(T_a^\top T_b) - 2\operatorname{tr}(T_a)\operatorname{tr}(T_b),
\label{eq:killing_metric}
\end{equation}

which is positive definite on $\mathfrak{sl}(K) = \ker(\operatorname{tr})$ and degenerate on the center $\mathbb{R} \cdot I$ (regularized by $\tilde{g} \to \tilde{g} + \epsilon I$ in practice). The natural gradient $\tilde{g}^{-1} \nabla_\phi \mathcal{F}$ has no free parameters: the metric is determined entirely by the algebra structure. This automatically handles the compact/non-compact asymmetry since the Killing form assigns different ``lengths'' to rotations versus shears.

\paragraph{4. Pullback natural gradient (position-dependent).}
The most geometrically principled approach uses the pullback of the Gram metric through the exponential map. The differential of $\exp$ at $\phi \in \mathfrak{g}$ in direction $T_a$ is:

\begin{equation}
d\exp_\phi(T_a) = \exp(\phi) \cdot \Psi(\mathrm{ad}_\phi)(T_a), \qquad \Psi(z) = \frac{e^z - 1}{z} = \sum_{k=0}^{\infty} \frac{z^k}{(k+1)!},
\end{equation}

where $\mathrm{ad}_\phi(X) = [\phi, X]$ is the adjoint action and $\Psi$ is the inverse Bernoulli function. The pullback metric on $\mathfrak{g}$ is:

\begin{equation}
\mathcal{G}_{ab}(\phi) = \left\langle \Psi(\mathrm{ad}_\phi)(T_a),\; \Psi(\mathrm{ad}_\phi)(T_b) \right\rangle_G,
\label{eq:pullback_metric}
\end{equation}

where $\langle X, Y \rangle_G = \operatorname{tr}(X^\top Y)$ is the Frobenius inner product. At $\phi = 0$, $\mathcal{G}_{ab}(0) = G_{ab}$ (the Gram matrix). For large $\|\phi\|$ in symmetric directions, $\mathcal{G}$ grows exponentially, automatically compensating the exponential amplification of $d\exp$ in the non-compact sector. The natural gradient $\mathcal{G}(\phi)^{-1} \nabla_\phi \mathcal{F}$ is position-dependent and exact (up to series truncation of $\Psi$), requiring the structure constants $f^c_{ab}$ defined by $[T_a, T_b] = \sum_c f^c_{ab} T_c$.

The computational cost of each mode scales as: clipping $O(n_{\mathrm{gen}})$, Cartan $O(n_{\mathrm{gen}}^2)$, Killing $O(n_{\mathrm{gen}}^2)$, pullback $O(n_{\mathrm{gen}}^3)$ (dominated by the linear solve $\mathcal{G}^{-1} g$). For the $\mathrm{GL}(10)$ model ($n_{\mathrm{gen}} = 100$), Killing form preconditioning adds negligible overhead. The pullback metric becomes advantageous when gauge frames develop large norms during training, as the position-dependent compensation prevents the gradient from overshooting in the symmetric directions where $\exp(\phi)$ amplifies exponentially.


\section{Variational Gradient Descent: Implementation and Numerical Methods}


\subsection{Natural Gradient Descent on the Gaussian Manifold}
Our implementation uses natural gradient descent for the statistical parameters, which exploits the information-geometric structure of the Gaussian manifold. Rather than treating the covariance matrices $\Sigma_i$ as Euclidean variables, we recognize them as points on the manifold of symmetric positive-definite matrices, equipped with the Fisher-Rao metric. The natural gradient at a point $\Sigma$ on this manifold is obtained by applying the inverse Fisher metric to the Euclidean gradient, yielding an intrinsic tangent vector that respects the manifold's curved geometry. Given Euclidean gradients $\nabla_\mu$ and $\nabla_\Sigma$, the natural gradient projections are
\begin{equation}
\delta\mu = -\Sigma\,\nabla_\mu, \qquad
\delta\Sigma = -2\,\Sigma\,\mathrm{sym}(\nabla_\Sigma)\,\Sigma,
\end{equation}
where $\mathrm{sym}(M) = \tfrac{1}{2}(M + M^\top)$.

The covariance update $\delta\Sigma$ is then passed to a manifold retraction (Section~\ref{sec:retraction}) that whitens it into the eigenbasis of $\Sigma$:
\begin{equation}
B = \Lambda^{-1/2}\, (V^\top \delta\Sigma\, V)\, \Lambda^{-1/2},
\end{equation}
where $\Sigma = V \Lambda V^\top$ is the eigendecomposition.  This whitening ensures that the resulting step is affine-invariant such that the same step size produces comparable effects regardless of the current scale of the covariance matrix.  Following the computation of the whitened tangent $B$, we apply a trust-region constraint by clipping its Frobenius norm to a maximum relative step size $\rho$. This prevents excessively large updates that could destabilize the manifold structure or lead to ill-conditioned covariances.

\subsection{Manifold Retraction for Covariance Matrices}\label{sec:retraction}
The manifold retraction is the process of mapping the tangent vector back to a valid point on the manifold. This is performed using the affine-invariant exponential map. Working in the eigenbasis of $\Sigma_k$ with eigenvalues $\lambda_j$, the whitened tangent $B$ is diagonalized $B = U\Lambda_B U^\top$ and the retraction is
\begin{equation}
\Sigma_{k+1} = V\,\Lambda^{1/2}\, U\, \exp(\tau \Lambda_B)\, U^\top\, \Lambda^{1/2}\, V^\top,
\end{equation}
where $\tau$ is the step size and the component-wise exponentials $\exp(\tau \lambda_B^{(j)})$ are clipped to $|\tau \lambda_B^{(j)}| \le 50$ to prevent floating-point overflow.  This procedure is provably SPD-preserving provided $\Sigma_k$ is SPD and $B$ is symmetric, which our implementation guarantees through explicit symmetrization at multiple stages.

Additional safeguards include post-retraction sanitization via
\texttt{sanitize\_sigma}, which symmetrizes the result, raises on any
floating-point anomalies (NaNs), and applies a spectral floor to the
eigenvalues to ensure they remain above a minimum threshold
$\epsilon_{\mathrm{SPD}} = 10^{-4}$, together with a condition-number cap
$\kappa_{\max} = 10^{4}$.  This spectral regularization is essential for
preventing numerical collapse of the covariance matrices during prolonged
gradient descent, particularly in regions where the free-energy landscape
becomes very flat.

\subsection{Gauge Frame Dynamics on $\mathrm{GL}(K)$}
For the gauge frame fields $\phi_i$, which live in the Lie algebra $\mathfrak{gl}(K)$ (or its subalgebra $\mathfrak{so}(3)$ in the fundamental representation), a different geometric structure must be respected. These fields parameterize the local gauge frames as $U_i = \exp(\sum_a \phi_i^a T_a)$, where $\{T_a\}$ are the Lie algebra generators, and the transport operators take the form $\Omega_{ij} = \exp(\phi_i)\exp(-\phi_j)$.  Gradients of the alignment energy $\beta_{ij}\,D_{\mathrm{KL}}(q_i \| \Omega_{ij}[q_j])$ with respect to $\phi_i$ are obtained via the chain rule through the transport operator: the KL gradient with respect to $\Omega$ is contracted against the differential $\partial\Omega_{ij}/\partial\phi_i^a$.

The differential of the matrix exponential relates infinitesimal variations $\delta \phi$ to variations in the group element via
\begin{equation}
\frac{\partial}{\partial \phi^a}\exp\!\Bigl(\textstyle\sum_b \phi^b T_b\Bigr) = \exp(X)\cdot Q_a,
\end{equation}
where $X = \sum_b \phi^b T_b$ and $Q_a = d\exp_X(T_a)$ is the derivative of the exponential map applied to the generator $T_a$.  For the $K\!=\!3$ fundamental representation of $\mathfrak{so}(3)$, we use the exact closed-form expression
\begin{equation}
Q_a = T_a - c_1(\theta)\,\mathrm{ad}_X(T_a) + c_2(\theta)\,\mathrm{ad}_X^2(T_a),
\end{equation}
where $\theta = \|\phi\|$, $c_1(\theta) = (1-\cos\theta)/\theta^2$, and $c_2(\theta) = (\theta - \sin\theta)/\theta^3$, with Taylor-expanded small-angle limits for $\theta < 10^{-4}$.  For higher-dimensional representations ($K > 3$), where the Cayley--Hamilton truncation requires $K$ terms, we instead compute $Q_a$ via the Fr\'{e}chet derivative of the matrix exponential using symmetric quadrature.

The resulting Euclidean gradient $\partial\mathcal{F}/\partial\phi^a$ can optionally be preconditioned using one of several Lie-algebraic metrics.  In the transformer implementation, these include: (i)~the \emph{Killing form} natural gradient, which uses the Cartan-involution-modified Killing form $\tilde{g}_{ab} = 2K\,\mathrm{tr}(T_a^\top T_b) - 2\,\mathrm{tr}(T_a)\mathrm{tr}(T_b)$ as a position-independent metric with no free parameters; (ii)~a \emph{Cartan decomposition} preconditioner that dampens the non-compact (symmetric) directions of $\mathfrak{gl}(K)$ relative to the compact $\mathfrak{so}(K)$ directions; and (iii)~the full \emph{pullback} natural gradient $G_{ab}(\phi) = \langle \Psi(\mathrm{ad}_X)(T_a),\, \Psi(\mathrm{ad}_X)(T_b) \rangle$, where $\Psi(z) = (e^z - 1)/z$, which captures the position-dependent curvature of the exponential map.

Once the gradient direction is determined, we apply a step with learning rate $\eta_\phi$ (typically $10^{-1}$), forming a candidate update $\phi_{\text{cand}} = \phi + \eta_\phi\,\delta \phi$. However, the Lie algebra has a natural periodic boundary (rotation angles wrap at $\pm \pi$) nd numerical drift can cause $\phi$ to wander outside the principal domain. To prevent this, we apply a retraction via \texttt{retract\_to\_principal\_ball}, which wraps angles modulo $2\pi$, applies an antipodal flip for angles exceeding $\pi$, and clamps the result to $\|\phi\| < \pi - \varepsilon_{\mathrm{margin}}$ with $\varepsilon_{\mathrm{margin}} = 10^{-2}$ to avoid numerical instability near the critical points where the exponential map degenerates.


\paragraph{Gradient validation.} All gauge frame gradient implementations are validated against finite-difference approximations $(\partial \mathcal{F}/\partial \phi_i^a \approx [\mathcal{F}(\phi_i + \epsilon e_a) - \mathcal{F}(\phi_i - \epsilon e_a)]/(2\epsilon))$ with relative error $< 10^{-5}$.


\bibliographystyle{plainnat}
\bibliography{references}

\end{document}
